\begin{afterword}
\summaryhead

The post answers readers who asked why the author wrote so much about physics. One purpose was to give a clear case where probability theory and Science disagree. The author holds that in physics simplicity can be compared strictly, by writing the rival theories as programs and counting lines of code, and the evidence is not in dispute. Many-worlds is simpler than collapse but makes no new experimental prediction.

So physicists who reject it are not breaking the rules of Science but obeying them. The ideal of Science says a new theory must win by experiment. The post grants that this is the ideal, not how science actually works. Teaching scientists Solomonoff induction would not fix this, because the rule that only experiment decides is what protects people from reasonable-sounding armchair theories; even Robert Aumann, a Bayesian theorist who believes in God, benefits from it. The reader must choose: reject the scientific method, or reject probability theory and believe in a single world, which the author calls insanity. A postscript warns that attempts to escape the dilemma will be answered in later posts.

\responsehead

The post is candid about its purpose: a clear case ``when it came time to break your allegiance to Science.'' It is also candid that its ``Science'' is an ideal, ``not how things do work in science.'' Its account of why science distrusts armchair reasoning is fair, and its Aumann example is accurate: Aumann let a test of the Bible codes decide, and reported that it failed to confirm them.

The dilemma sets probability theory against that narrow ideal. The rule that a new theory must win by novel prediction is one side of a long debate: it has been attributed to Bacon and Descartes and was defended by Whewell, while Mill and Keynes rejected it. And physicists have chosen between empirically equivalent theories on grounds of simplicity: special relativity is preferred to Lorentz's ether theory for its lack of an undetectable ether. That case supports the author's view that simplicity can decide. It also shows that preferring the simpler theory need not mean rejecting the scientific method.

The other horn claims more than probability theory gives. The ``strict'' comparison by lines of code is not carried out, here or earlier, and what the many-worlds program must contain is disputed: a rule for probabilities, and perhaps the observer's location. Probability theory does not choose the prior on which the comparison depends. The rival offered is the textbook collapse postulate, not the specified collapse theories, and Bohm's theory, one world with no collapse, is not mentioned. So a Bayesian can doubt many-worlds without rejecting probability theory. These points are made in full in the notes on ``Decoherence is Simple'' and ``Many Worlds, One Best Guess.''

The post handles doubt in the same way. The commenter it quotes doubted a conclusion learned from one advocate; the post recasts the doubt as the pull of Science's rules. A single world is called ``insanity,'' more than the author's own hedge two days earlier allowed. And objections are called ``clever ways to wriggle out'' before any is made.

In short: a candid statement of the aim to turn readers from Science to Bayes, which sets probability theory against a narrow ideal of Science, treats a disputed simplicity argument as probability theory itself, and leaves out the positions between its two horns.
\end{afterword}
